% Options for packages loaded elsewhere
% Options for packages loaded elsewhere
\PassOptionsToPackage{unicode}{hyperref}
\PassOptionsToPackage{hyphens}{url}
\PassOptionsToPackage{dvipsnames,svgnames,x11names}{xcolor}
%
\documentclass[
  letterpaper,
  DIV=11,
  numbers=noendperiod]{scrartcl}
\usepackage{xcolor}
\usepackage{amsmath,amssymb}
\setcounter{secnumdepth}{-\maxdimen} % remove section numbering
\usepackage{iftex}
\ifPDFTeX
  \usepackage[T1]{fontenc}
  \usepackage[utf8]{inputenc}
  \usepackage{textcomp} % provide euro and other symbols
\else % if luatex or xetex
  \usepackage{unicode-math} % this also loads fontspec
  \defaultfontfeatures{Scale=MatchLowercase}
  \defaultfontfeatures[\rmfamily]{Ligatures=TeX,Scale=1}
\fi
\usepackage{lmodern}
\ifPDFTeX\else
  % xetex/luatex font selection
\fi
% Use upquote if available, for straight quotes in verbatim environments
\IfFileExists{upquote.sty}{\usepackage{upquote}}{}
\IfFileExists{microtype.sty}{% use microtype if available
  \usepackage[]{microtype}
  \UseMicrotypeSet[protrusion]{basicmath} % disable protrusion for tt fonts
}{}
\makeatletter
\@ifundefined{KOMAClassName}{% if non-KOMA class
  \IfFileExists{parskip.sty}{%
    \usepackage{parskip}
  }{% else
    \setlength{\parindent}{0pt}
    \setlength{\parskip}{6pt plus 2pt minus 1pt}}
}{% if KOMA class
  \KOMAoptions{parskip=half}}
\makeatother
% Make \paragraph and \subparagraph free-standing
\makeatletter
\ifx\paragraph\undefined\else
  \let\oldparagraph\paragraph
  \renewcommand{\paragraph}{
    \@ifstar
      \xxxParagraphStar
      \xxxParagraphNoStar
  }
  \newcommand{\xxxParagraphStar}[1]{\oldparagraph*{#1}\mbox{}}
  \newcommand{\xxxParagraphNoStar}[1]{\oldparagraph{#1}\mbox{}}
\fi
\ifx\subparagraph\undefined\else
  \let\oldsubparagraph\subparagraph
  \renewcommand{\subparagraph}{
    \@ifstar
      \xxxSubParagraphStar
      \xxxSubParagraphNoStar
  }
  \newcommand{\xxxSubParagraphStar}[1]{\oldsubparagraph*{#1}\mbox{}}
  \newcommand{\xxxSubParagraphNoStar}[1]{\oldsubparagraph{#1}\mbox{}}
\fi
\makeatother

\usepackage{color}
\usepackage{fancyvrb}
\newcommand{\VerbBar}{|}
\newcommand{\VERB}{\Verb[commandchars=\\\{\}]}
\DefineVerbatimEnvironment{Highlighting}{Verbatim}{commandchars=\\\{\}}
% Add ',fontsize=\small' for more characters per line
\usepackage{framed}
\definecolor{shadecolor}{RGB}{241,243,245}
\newenvironment{Shaded}{\begin{snugshade}}{\end{snugshade}}
\newcommand{\AlertTok}[1]{\textcolor[rgb]{0.68,0.00,0.00}{#1}}
\newcommand{\AnnotationTok}[1]{\textcolor[rgb]{0.37,0.37,0.37}{#1}}
\newcommand{\AttributeTok}[1]{\textcolor[rgb]{0.40,0.45,0.13}{#1}}
\newcommand{\BaseNTok}[1]{\textcolor[rgb]{0.68,0.00,0.00}{#1}}
\newcommand{\BuiltInTok}[1]{\textcolor[rgb]{0.00,0.23,0.31}{#1}}
\newcommand{\CharTok}[1]{\textcolor[rgb]{0.13,0.47,0.30}{#1}}
\newcommand{\CommentTok}[1]{\textcolor[rgb]{0.37,0.37,0.37}{#1}}
\newcommand{\CommentVarTok}[1]{\textcolor[rgb]{0.37,0.37,0.37}{\textit{#1}}}
\newcommand{\ConstantTok}[1]{\textcolor[rgb]{0.56,0.35,0.01}{#1}}
\newcommand{\ControlFlowTok}[1]{\textcolor[rgb]{0.00,0.23,0.31}{\textbf{#1}}}
\newcommand{\DataTypeTok}[1]{\textcolor[rgb]{0.68,0.00,0.00}{#1}}
\newcommand{\DecValTok}[1]{\textcolor[rgb]{0.68,0.00,0.00}{#1}}
\newcommand{\DocumentationTok}[1]{\textcolor[rgb]{0.37,0.37,0.37}{\textit{#1}}}
\newcommand{\ErrorTok}[1]{\textcolor[rgb]{0.68,0.00,0.00}{#1}}
\newcommand{\ExtensionTok}[1]{\textcolor[rgb]{0.00,0.23,0.31}{#1}}
\newcommand{\FloatTok}[1]{\textcolor[rgb]{0.68,0.00,0.00}{#1}}
\newcommand{\FunctionTok}[1]{\textcolor[rgb]{0.28,0.35,0.67}{#1}}
\newcommand{\ImportTok}[1]{\textcolor[rgb]{0.00,0.46,0.62}{#1}}
\newcommand{\InformationTok}[1]{\textcolor[rgb]{0.37,0.37,0.37}{#1}}
\newcommand{\KeywordTok}[1]{\textcolor[rgb]{0.00,0.23,0.31}{\textbf{#1}}}
\newcommand{\NormalTok}[1]{\textcolor[rgb]{0.00,0.23,0.31}{#1}}
\newcommand{\OperatorTok}[1]{\textcolor[rgb]{0.37,0.37,0.37}{#1}}
\newcommand{\OtherTok}[1]{\textcolor[rgb]{0.00,0.23,0.31}{#1}}
\newcommand{\PreprocessorTok}[1]{\textcolor[rgb]{0.68,0.00,0.00}{#1}}
\newcommand{\RegionMarkerTok}[1]{\textcolor[rgb]{0.00,0.23,0.31}{#1}}
\newcommand{\SpecialCharTok}[1]{\textcolor[rgb]{0.37,0.37,0.37}{#1}}
\newcommand{\SpecialStringTok}[1]{\textcolor[rgb]{0.13,0.47,0.30}{#1}}
\newcommand{\StringTok}[1]{\textcolor[rgb]{0.13,0.47,0.30}{#1}}
\newcommand{\VariableTok}[1]{\textcolor[rgb]{0.07,0.07,0.07}{#1}}
\newcommand{\VerbatimStringTok}[1]{\textcolor[rgb]{0.13,0.47,0.30}{#1}}
\newcommand{\WarningTok}[1]{\textcolor[rgb]{0.37,0.37,0.37}{\textit{#1}}}

\usepackage{longtable,booktabs,array}
\newcounter{none} % for unnumbered tables
\usepackage{calc} % for calculating minipage widths
% Correct order of tables after \paragraph or \subparagraph
\usepackage{etoolbox}
\makeatletter
\patchcmd\longtable{\par}{\if@noskipsec\mbox{}\fi\par}{}{}
\makeatother
% Allow footnotes in longtable head/foot
\IfFileExists{footnotehyper.sty}{\usepackage{footnotehyper}}{\usepackage{footnote}}
\makesavenoteenv{longtable}
\usepackage{graphicx}
\makeatletter
\newsavebox\pandoc@box
\newcommand*\pandocbounded[1]{% scales image to fit in text height/width
  \sbox\pandoc@box{#1}%
  \Gscale@div\@tempa{\textheight}{\dimexpr\ht\pandoc@box+\dp\pandoc@box\relax}%
  \Gscale@div\@tempb{\linewidth}{\wd\pandoc@box}%
  \ifdim\@tempb\p@<\@tempa\p@\let\@tempa\@tempb\fi% select the smaller of both
  \ifdim\@tempa\p@<\p@\scalebox{\@tempa}{\usebox\pandoc@box}%
  \else\usebox{\pandoc@box}%
  \fi%
}
% Set default figure placement to htbp
\def\fps@figure{htbp}
\makeatother


% definitions for citeproc citations
\NewDocumentCommand\citeproctext{}{}
\NewDocumentCommand\citeproc{mm}{%
  \begingroup\def\citeproctext{#2}\cite{#1}\endgroup}
\makeatletter
 % allow citations to break across lines
 \let\@cite@ofmt\@firstofone
 % avoid brackets around text for \cite:
 \def\@biblabel#1{}
 \def\@cite#1#2{{#1\if@tempswa , #2\fi}}
\makeatother
\newlength{\cslhangindent}
\setlength{\cslhangindent}{1.5em}
\newlength{\csllabelwidth}
\setlength{\csllabelwidth}{3em}
\newenvironment{CSLReferences}[2] % #1 hanging-indent, #2 entry-spacing
 {\begin{list}{}{%
  \setlength{\itemindent}{0pt}
  \setlength{\leftmargin}{0pt}
  \setlength{\parsep}{0pt}
  % turn on hanging indent if param 1 is 1
  \ifodd #1
   \setlength{\leftmargin}{\cslhangindent}
   \setlength{\itemindent}{-1\cslhangindent}
  \fi
  % set entry spacing
  \setlength{\itemsep}{#2\baselineskip}}}
 {\end{list}}
\usepackage{calc}
\newcommand{\CSLBlock}[1]{\hfill\break\parbox[t]{\linewidth}{\strut\ignorespaces#1\strut}}
\newcommand{\CSLLeftMargin}[1]{\parbox[t]{\csllabelwidth}{\strut#1\strut}}
\newcommand{\CSLRightInline}[1]{\parbox[t]{\linewidth - \csllabelwidth}{\strut#1\strut}}
\newcommand{\CSLIndent}[1]{\hspace{\cslhangindent}#1}



\setlength{\emergencystretch}{3em} % prevent overfull lines

\providecommand{\tightlist}{%
  \setlength{\itemsep}{0pt}\setlength{\parskip}{0pt}}



 


\KOMAoption{captions}{tableheading}
\makeatletter
\@ifpackageloaded{caption}{}{\usepackage{caption}}
\AtBeginDocument{%
\ifdefined\contentsname
  \renewcommand*\contentsname{Table of contents}
\else
  \newcommand\contentsname{Table of contents}
\fi
\ifdefined\listfigurename
  \renewcommand*\listfigurename{List of Figures}
\else
  \newcommand\listfigurename{List of Figures}
\fi
\ifdefined\listtablename
  \renewcommand*\listtablename{List of Tables}
\else
  \newcommand\listtablename{List of Tables}
\fi
\ifdefined\figurename
  \renewcommand*\figurename{Figure}
\else
  \newcommand\figurename{Figure}
\fi
\ifdefined\tablename
  \renewcommand*\tablename{Table}
\else
  \newcommand\tablename{Table}
\fi
}
\@ifpackageloaded{float}{}{\usepackage{float}}
\floatstyle{ruled}
\@ifundefined{c@chapter}{\newfloat{codelisting}{h}{lop}}{\newfloat{codelisting}{h}{lop}[chapter]}
\floatname{codelisting}{Listing}
\newcommand*\listoflistings{\listof{codelisting}{List of Listings}}
\makeatother
\makeatletter
\makeatother
\makeatletter
\@ifpackageloaded{caption}{}{\usepackage{caption}}
\@ifpackageloaded{subcaption}{}{\usepackage{subcaption}}
\makeatother
\usepackage{bookmark}
\IfFileExists{xurl.sty}{\usepackage{xurl}}{} % add URL line breaks if available
\urlstyle{same}
\hypersetup{
  colorlinks=true,
  linkcolor={blue},
  filecolor={Maroon},
  citecolor={Blue},
  urlcolor={Blue},
  pdfcreator={LaTeX via pandoc}}


\author{}
\date{}
\begin{document}


\subsection{}\label{section}

\textbf{Neural Hydrology}

Deep Learning for Streamflow Prediction under Climate Change

Author: Yuxin Qiu

Supervisor: Henri Funk

Seminar: Climate Change Statistics

University: LMU Munich

Date: 2026-06-09  ·  Interim Meeting

\begin{center}\rule{0.5\linewidth}{0.5pt}\end{center}

\subsection{Contents}\label{contents}

\begin{itemize}
\tightlist
\item
  Research theory base.
\item
  Project question and hypotheses.
\item
  Data, model setup, and evaluation criteria.
\item
  Group-level empirical results.
\item
  Basin-level empirical results.
\item
  Robustness checks and interpretation.
\item
  Practical guidance, limitations, and next steps.
\end{itemize}

\begin{center}\rule{0.5\linewidth}{0.5pt}\end{center}

\subsection{What is Neural Hydrology?}\label{what-is-neural-hydrology}

\textbf{Streamflow} is the volume of water flowing through a river
channel at a given point per unit time --- typically measured in mm/day
or m³/s. It is the integrated output of the entire catchment:
precipitation, evaporation, soil infiltration, and groundwater dynamics
all feed into it (Addor et al. 2017).

\textbf{Rainfall-runoff modelling} is the task of predicting daily
streamflow \(Q\) from meteorological inputs (precipitation, temperature,
radiation) --- a long-standing central problem in hydrology (Kratzert et
al. 2018).

. . .

\textbf{Neural Hydrology} is the application of deep learning to this
problem. Instead of hand-crafting physical equations, a neural network
learns the input-to-streamflow mapping directly from data, across
hundreds of catchments simultaneously.

\subsection{Milestones in Neural
Hydrology}\label{milestones-in-neural-hydrology}

\begin{itemize}
\tightlist
\item
  \textbf{2018} --- First successful LSTM rainfall--runoff model,
  outperforming the calibrated SAC-SMA + Snow17 benchmark (Kratzert et
  al. 2018)
\item
  \textbf{2019} --- EA-LSTM: one model trained on 531 CAMELS basins
  simultaneously (Kratzert et al. 2019)
\item
  \textbf{2021} --- MC-LSTM: mass conservation enforced as a structural
  guarantee (Hoedt et al. 2021)
\item
  \textbf{2022} --- NeuralHydrology released as the first unified
  open-source deep-learning framework for hydrology (Kratzert et al.
  2022)
\item
  \textbf{2023} --- Caravan: global large-sample hydrology benchmark
  with 6,830 catchments (Kratzert et al. 2023)
\end{itemize}

\begin{center}\rule{0.5\linewidth}{0.5pt}\end{center}

\subsection{Why Streamflow Prediction
Matters}\label{why-streamflow-prediction-matters}

\begin{itemize}
\tightlist
\item
  \textbf{Floods} are the most frequent natural disaster --- 57 million
  people affected in 2022 alone (Centre for Research on the Epidemiology
  of Disasters (CRED) 2022)
\item
  \textbf{Floods} are among the most frequent natural disasters and
  create large annual human and economic impacts (Centre for Research on
  the Epidemiology of Disasters (CRED) 2022)
\item
  Accurate \textbf{streamflow forecasts} are critical for flood
  early-warning, reservoir management, and hydropower scheduling
\item
  Under \textbf{climate change}, precipitation extremes intensify,
  seasonal patterns shift, and snowmelt regimes change (IPCC 2021)
\item
  Hydrological models must be \textbf{robust to non-stationarity} ---
  conditions outside the historical training distribution
\end{itemize}

. . .

\begin{quote}
Milly et al.~(2008) framed the challenge as \emph{``Stationarity is
dead: Whither water management?''} (Milly et al. 2008)
\end{quote}

\begin{center}\rule{0.5\linewidth}{0.5pt}\end{center}

\subsection{The Deep Learning
Breakthrough}\label{the-deep-learning-breakthrough}

\begin{center}
\pandocbounded{\includegraphics[keepaspectratio]{slides_files/figure-pdf/model-comparison-bar-1.pdf}}
\end{center}

Process-based models (SAC-SMA+Snow-17, and others) require calibrating
10--40 parameters per catchment --- not scalable to ungauged basins
(Kratzert et al. 2018). Individual LSTMs were trained separately for
each of the \textbf{241 basins} with no manual calibration; mean NSE =
\textbf{0.63} versus SAC-SMA mean NSE = \textbf{0.58} (Kratzert et al.
2018).

\begin{center}\rule{0.5\linewidth}{0.5pt}\end{center}

\subsection{LSTM and EA-LSTM}\label{lstm-and-ea-lstm}

The LSTM (Hochreiter and Schmidhuber 1997) maintains a \textbf{cell
state} \(\mathbf{c}_t\) as hydrological memory:

\[\mathbf{c}_t = \underbrace{\mathbf{f}_t \odot \mathbf{c}_{t-1}}_{\text{long-term memory}} + \underbrace{\mathbf{i}_t \odot \tanh(\mathbf{W}_g [\mathbf{h}_{t-1}, \mathbf{x}_t] + \mathbf{b}_g)}_{\text{new information}} \qquad \mathbf{h}_t = \mathbf{o}_t \odot \tanh(\mathbf{c}_t)\]

Cell state \(\approx\) soil moisture / groundwater · Forget gate
\(\approx\) drainage · Input gate \(\approx\) infiltration

. . .

\textbf{Entity-Aware LSTM} (Kratzert et al. 2019) --- conditions the
input gate on static catchment attributes \(\mathbf{x}_s\):

\[\mathbf{i} = \sigma(\mathbf{W}_i \mathbf{x}_s + \mathbf{b}_i)\]

One model trained simultaneously on \textbf{531 CAMELS-US basins}:
median NSE \textbf{0.72} --- outperforming the best calibrated process
model (SAC-SMA, NSE 0.58) by a wide margin.

The \textbf{Nash-Sutcliffe Efficiency}
\(\text{NSE} = 1 - \frac{\sum(\hat{q}_t - q_t)^2}{\sum(\bar{q} - q_t)^2}\)
ranges from \(-\infty\) to 1. In several hydrological studies, NSE
values above 0.6 are treated as reasonably good, but this remains
context dependent.

\begin{center}\rule{0.5\linewidth}{0.5pt}\end{center}

\subsection{NeuralHydrology: Python
Library}\label{neuralhydrology-python-library}

\textbf{NeuralHydrology} (Kratzert et al. 2022) --- open-source PyTorch
library built for reproducible deep learning in hydrology:

\begin{itemize}
\tightlist
\item
  \textbf{Config-file driven} --- a single \texttt{.yml} file specifies
  the full experiment; no source code changes between runs
\item
  \textbf{Model zoo} --- 15+ architectures: CudaLSTM, EA-LSTM, MC-LSTM,
  MTS-LSTM, Transformer, Mamba, xLSTM, HybridModel
\item
  \textbf{Dataset zoo} --- CAMELS-US/GB/DE/AUS/CL, Caravan
  (\textgreater6 000 basins globally)
\item
  \textbf{Probabilistic heads} --- GMM, CMAL, UMAL for uncertainty
  quantification
\item
  \textbf{Evaluation} --- NSE, KGE, MSE, hydrological signatures
  built-in
\end{itemize}

\begin{Shaded}
\begin{Highlighting}[]
\ExtensionTok{uv}\NormalTok{ pip install neuralhydrology}
\ExtensionTok{nh{-}run}\NormalTok{ train    }\AttributeTok{{-}{-}config{-}file}\NormalTok{ config.yml}
\ExtensionTok{nh{-}run}\NormalTok{ evaluate }\AttributeTok{{-}{-}run{-}dir}\NormalTok{    runs/my\_run/}
\end{Highlighting}
\end{Shaded}

\begin{center}\rule{0.5\linewidth}{0.5pt}\end{center}

\subsection{Formal Question and Testable
Hypotheses}\label{formal-question-and-testable-hypotheses}

\begin{quote}
Can global pretraining + local fine-tuning outperform local-from-scratch
training for regional streamflow prediction?
\end{quote}

\textbf{Question}: Does transfer learning provide stable gains across
heterogeneous basin groups?

\textbf{H1}: Fine-tune NSE \textgreater{} Local NSE for most groups.

\textbf{H2}: Fine-tune NSE \textgreater{} Global NSE for most groups.

\textbf{Expected ordering}: \[
    ext{NSE}_{\text{Fine-tune}} > \text{NSE}_{\text{Global}} > \text{NSE}_{\text{Local}}
\]

\begin{center}\rule{0.5\linewidth}{0.5pt}\end{center}

\subsection{Why Three-Way Comparison Is
Necessary}\label{why-three-way-comparison-is-necessary}

\begin{itemize}
\tightlist
\item
  Local baseline answers: how much can target-region data alone learn?
\item
  Global baseline answers: how much transfer is possible without
  adaptation?
\item
  Fine-tune model answers: how much additional value comes from
  adaptation?
\end{itemize}

Without all three, transfer benefit cannot be separated from model
capacity.

\begin{center}\rule{0.5\linewidth}{0.5pt}\end{center}

\subsection{Experimental Arms and
Fairness}\label{experimental-arms-and-fairness}

{\def\LTcaptype{none} % do not increment counter
\begin{longtable}[]{@{}
  >{\raggedright\arraybackslash}p{(\linewidth - 6\tabcolsep) * \real{0.2317}}
  >{\raggedright\arraybackslash}p{(\linewidth - 6\tabcolsep) * \real{0.2195}}
  >{\centering\arraybackslash}p{(\linewidth - 6\tabcolsep) * \real{0.2073}}
  >{\raggedright\arraybackslash}p{(\linewidth - 6\tabcolsep) * \real{0.3415}}@{}}
\toprule\noalign{}
\begin{minipage}[b]{\linewidth}\raggedright
Arm
\end{minipage} & \begin{minipage}[b]{\linewidth}\raggedright
Initialization
\end{minipage} & \begin{minipage}[b]{\linewidth}\centering
Train\_on\_Target
\end{minipage} & \begin{minipage}[b]{\linewidth}\raggedright
Primary\_Purpose
\end{minipage} \\
\midrule\noalign{}
\endhead
\bottomrule\noalign{}
\endlastfoot
Local & Random & Yes & Local baseline \\
Global & Global pretrained & No & Transfer without adaptation \\
Global + Fine-tune & Global pretrained & Yes & Transfer with
adaptation \\
\end{longtable}
}

\begin{center}\rule{0.5\linewidth}{0.5pt}\end{center}

\subsection{Model and Training Setup}\label{model-and-training-setup}

{\def\LTcaptype{none} % do not increment counter
\begin{longtable}[]{@{}ll@{}}
\toprule\noalign{}
Item & Value \\
\midrule\noalign{}
\endhead
\bottomrule\noalign{}
\endlastfoot
Architecture & LSTM-based NH setup \\
Global pretraining scope & 531 CAMELS-US basins \\
Groups & 18 Daymet-based groups \\
Random seed & 1001 \\
Evaluation metrics & NSE (primary), KGE (support) \\
\end{longtable}
}

\begin{center}\rule{0.5\linewidth}{0.5pt}\end{center}

\subsection{Data and Result Files
Used}\label{data-and-result-files-used}

\begin{itemize}
\tightlist
\item
  Group summary:
  ../rq3\_finetune/results\_folder\_groups\_summary\_with\_global\_true.csv
\item
  Basin-level detail:
  ../rq3\_finetune/results\_folder\_groups\_all\_basins.csv
\item
  Per-group local vs fine-tune details:
  results\_folder\_XX\_local\_vs\_finetune.csv
\item
  Per-group global evaluation details:
  results\_folder\_XX\_global\_eval\_true.csv
\end{itemize}

\begin{center}\rule{0.5\linewidth}{0.5pt}\end{center}

\subsection{Data Preparation for RQ3
Slides}\label{data-preparation-for-rq3-slides}

\begin{center}\rule{0.5\linewidth}{0.5pt}\end{center}

\subsection{Coverage Across Groups}\label{coverage-across-groups}

\begin{center}
\pandocbounded{\includegraphics[keepaspectratio]{slides_files/figure-pdf/rq3-coverage-bars-1.pdf}}
\end{center}

\begin{center}\rule{0.5\linewidth}{0.5pt}\end{center}

\subsection{Group Size Distribution}\label{group-size-distribution}

\begin{center}
\pandocbounded{\includegraphics[keepaspectratio]{slides_files/figure-pdf/rq3-group-size-dist-1.pdf}}
\end{center}

\begin{longtable}[]{@{}ll@{}}
\caption{Distribution Summary (n=18 groups)}\tabularnewline
\toprule\noalign{}
Metric & Value \\
\midrule\noalign{}
\endfirsthead
\toprule\noalign{}
Metric & Value \\
\midrule\noalign{}
\endhead
\bottomrule\noalign{}
\endlastfoot
Min & 2 \\
Q1 & 14 \\
Median & 26 \\
Q3 & 34 \\
Max & 79 \\
Mean & 29.6 \\
SD & 23.4 \\
\end{longtable}

\begin{center}\rule{0.5\linewidth}{0.5pt}\end{center}

\subsection{Overall Weighted NSE
Comparison}\label{overall-weighted-nse-comparison}

\begin{center}
\pandocbounded{\includegraphics[keepaspectratio]{slides_files/figure-pdf/rq3-overall-nse-1.pdf}}
\end{center}

\begin{center}\rule{0.5\linewidth}{0.5pt}\end{center}

\subsection{Overall Weighted KGE
Comparison}\label{overall-weighted-kge-comparison}

\begin{center}
\pandocbounded{\includegraphics[keepaspectratio]{slides_files/figure-pdf/rq3-overall-kge-1.pdf}}
\end{center}

\begin{center}\rule{0.5\linewidth}{0.5pt}\end{center}

\subsection{Group-Level NSE: Three-Way
Comparison}\label{group-level-nse-three-way-comparison}

\begin{center}
\pandocbounded{\includegraphics[keepaspectratio]{slides_files/figure-pdf/rq3-three-way-by-group-1.pdf}}
\end{center}

\begin{center}\rule{0.5\linewidth}{0.5pt}\end{center}

\subsection{Group-Level Delta NSE (FT vs Local and FT vs
Global)}\label{group-level-delta-nse-ft-vs-local-and-ft-vs-global}

\begin{center}
\pandocbounded{\includegraphics[keepaspectratio]{slides_files/figure-pdf/rq3-group-delta-nse-1.pdf}}
\end{center}

\begin{center}\rule{0.5\linewidth}{0.5pt}\end{center}

\subsection{Group-Level Delta KGE (FT vs
Local)}\label{group-level-delta-kge-ft-vs-local}

\begin{center}
\pandocbounded{\includegraphics[keepaspectratio]{slides_files/figure-pdf/rq3-group-delta-kge-1.pdf}}
\end{center}

\begin{center}\rule{0.5\linewidth}{0.5pt}\end{center}

\subsection{Group-Level Improvement Rate (Within
Group)}\label{group-level-improvement-rate-within-group}

\begin{center}
\pandocbounded{\includegraphics[keepaspectratio]{slides_files/figure-pdf/rq3-improve-rate-1.pdf}}
\end{center}

\begin{center}\rule{0.5\linewidth}{0.5pt}\end{center}

\subsection{Local Baseline vs Fine-Tuned Performance (Group
Means)}\label{local-baseline-vs-fine-tuned-performance-group-means}

\begin{center}
\pandocbounded{\includegraphics[keepaspectratio]{slides_files/figure-pdf/rq3-scatter-local-ft-1.pdf}}
\end{center}

\begin{center}\rule{0.5\linewidth}{0.5pt}\end{center}

\subsection{Global Baseline vs Fine-Tuned Performance (Group
Means)}\label{global-baseline-vs-fine-tuned-performance-group-means}

\begin{center}
\pandocbounded{\includegraphics[keepaspectratio]{slides_files/figure-pdf/rq3-scatter-global-ft-1.pdf}}
\end{center}

\begin{center}\rule{0.5\linewidth}{0.5pt}\end{center}

\subsection{Basin-Level NSE
Distribution}\label{basin-level-nse-distribution}

\begin{center}
\pandocbounded{\includegraphics[keepaspectratio]{slides_files/figure-pdf/rq3-basin-boxplot-1.pdf}}
\end{center}

\begin{center}\rule{0.5\linewidth}{0.5pt}\end{center}

\subsection{Basin-Level Delta NSE
Distribution}\label{basin-level-delta-nse-distribution}

\begin{center}
\pandocbounded{\includegraphics[keepaspectratio]{slides_files/figure-pdf/rq3-basin-delta-hist-1.pdf}}
\end{center}

\begin{center}\rule{0.5\linewidth}{0.5pt}\end{center}

\subsection{Basin-Level ECDF of NSE}\label{basin-level-ecdf-of-nse}

\begin{center}
\pandocbounded{\includegraphics[keepaspectratio]{slides_files/figure-pdf/rq3-basin-ecdf-1.pdf}}
\end{center}

\begin{center}\rule{0.5\linewidth}{0.5pt}\end{center}

\subsection{Top Basin Improvements
(NSE)}\label{top-basin-improvements-nse}

{\def\LTcaptype{none} % do not increment counter
\begin{longtable}[]{@{}llcccc@{}}
\toprule\noalign{}
& basin & group & local\_nse & ft\_nse & delta\_nse \\
\midrule\noalign{}
\endhead
\bottomrule\noalign{}
\endlastfoot
428 & 9510200 & 15 & 0.122 & 0.765 & 0.643 \\
346 & 7145700 & 11 & 0.107 & 0.729 & 0.622 \\
363 & 8023080 & 12 & 0.206 & 0.814 & 0.608 \\
429 & 9512280 & 15 & 0.090 & 0.691 & 0.601 \\
334 & 6910800 & 10 & 0.205 & 0.804 & 0.599 \\
379 & 8164300 & 12 & 0.219 & 0.816 & 0.597 \\
335 & 6911900 & 10 & 0.181 & 0.769 & 0.589 \\
513 & 11151300 & 18 & 0.314 & 0.888 & 0.574 \\
434 & 10336660 & 16 & 0.146 & 0.717 & 0.571 \\
433 & 10336645 & 16 & 0.277 & 0.847 & 0.570 \\
355 & 7263295 & 11 & 0.218 & 0.787 & 0.569 \\
242 & 3463300 & 06 & 0.160 & 0.726 & 0.566 \\
\end{longtable}
}

\begin{center}\rule{0.5\linewidth}{0.5pt}\end{center}

\subsection{Basins with Negative Delta
NSE}\label{basins-with-negative-delta-nse}

{\def\LTcaptype{none} % do not increment counter
\begin{longtable}[]{@{}llcccc@{}}
\toprule\noalign{}
& basin & group & local\_nse & ft\_nse & delta\_nse \\
\midrule\noalign{}
\endhead
\bottomrule\noalign{}
\endlastfoot
421 & 9484600 & 15 & 0.163 & -0.400 & -0.562 \\
464 & 12383500 & 17 & 0.604 & 0.342 & -0.262 \\
369 & 8082700 & 12 & 0.136 & -0.063 & -0.199 \\
463 & 12381400 & 17 & 0.712 & 0.571 & -0.141 \\
95 & 2046000 & 03 & 0.661 & 0.536 & -0.125 \\
443 & 12048000 & 17 & 0.232 & 0.138 & -0.094 \\
509 & 11124500 & 18 & 0.569 & 0.479 & -0.090 \\
306 & 6409000 & 10 & 0.125 & 0.075 & -0.050 \\
461 & 12375900 & 17 & 0.799 & 0.762 & -0.037 \\
43 & 1484100 & 02 & 0.603 & 0.567 & -0.037 \\
117 & 2112360 & 03 & 0.541 & 0.506 & -0.035 \\
465 & 12390700 & 17 & 0.656 & 0.634 & -0.023 \\
\end{longtable}
}

\begin{center}\rule{0.5\linewidth}{0.5pt}\end{center}

\subsection{Sample Size vs Transfer
Gain}\label{sample-size-vs-transfer-gain}

\begin{center}
\pandocbounded{\includegraphics[keepaspectratio]{slides_files/figure-pdf/rq3-size-vs-gain-1.pdf}}
\end{center}

\begin{center}\rule{0.5\linewidth}{0.5pt}\end{center}

\subsection{Local Baseline Quality vs Transfer
Gain}\label{local-baseline-quality-vs-transfer-gain}

\begin{center}
\pandocbounded{\includegraphics[keepaspectratio]{slides_files/figure-pdf/rq3-local-vs-gain-1.pdf}}
\end{center}

\begin{center}\rule{0.5\linewidth}{0.5pt}\end{center}

\subsection{Global Baseline Quality vs FT-over-Global
Gain}\label{global-baseline-quality-vs-ft-over-global-gain}

\begin{center}
\pandocbounded{\includegraphics[keepaspectratio]{slides_files/figure-pdf/rq3-global-vs-gain-1.pdf}}
\end{center}

\begin{center}\rule{0.5\linewidth}{0.5pt}\end{center}

\subsection{Group Ranking Table (NSE)}\label{group-ranking-table-nse}

{\def\LTcaptype{none} % do not increment counter
\begin{longtable}[]{@{}
  >{\raggedright\arraybackslash}p{(\linewidth - 14\tabcolsep) * \real{0.0265}}
  >{\centering\arraybackslash}p{(\linewidth - 14\tabcolsep) * \real{0.0619}}
  >{\centering\arraybackslash}p{(\linewidth - 14\tabcolsep) * \real{0.0885}}
  >{\centering\arraybackslash}p{(\linewidth - 14\tabcolsep) * \real{0.1416}}
  >{\centering\arraybackslash}p{(\linewidth - 14\tabcolsep) * \real{0.1504}}
  >{\centering\arraybackslash}p{(\linewidth - 14\tabcolsep) * \real{0.1150}}
  >{\centering\arraybackslash}p{(\linewidth - 14\tabcolsep) * \real{0.2035}}
  >{\centering\arraybackslash}p{(\linewidth - 14\tabcolsep) * \real{0.2124}}@{}}
\toprule\noalign{}
\begin{minipage}[b]{\linewidth}\raggedright
\end{minipage} & \begin{minipage}[b]{\linewidth}\centering
group
\end{minipage} & \begin{minipage}[b]{\linewidth}\centering
n\_basins
\end{minipage} & \begin{minipage}[b]{\linewidth}\centering
local\_nse\_mean
\end{minipage} & \begin{minipage}[b]{\linewidth}\centering
global\_nse\_mean
\end{minipage} & \begin{minipage}[b]{\linewidth}\centering
ft\_nse\_mean
\end{minipage} & \begin{minipage}[b]{\linewidth}\centering
delta\_ft\_vs\_local\_nse
\end{minipage} & \begin{minipage}[b]{\linewidth}\centering
delta\_ft\_vs\_global\_nse
\end{minipage} \\
\midrule\noalign{}
\endhead
\bottomrule\noalign{}
\endlastfoot
14 & 14 & 15 & 0.655 & 0.803 & 0.834 & 0.178 & 0.030 \\
17 & 17 & 72 & 0.680 & 0.772 & 0.795 & 0.115 & 0.023 \\
16 & 16 & 5 & 0.353 & 0.778 & 0.785 & 0.431 & 0.007 \\
1 & 01 & 25 & 0.520 & 0.725 & 0.769 & 0.249 & 0.044 \\
6 & 06 & 16 & 0.315 & 0.732 & 0.768 & 0.454 & 0.037 \\
18 & 18 & 26 & 0.469 & 0.740 & 0.766 & 0.298 & 0.026 \\
5 & 05 & 35 & 0.527 & 0.704 & 0.731 & 0.204 & 0.027 \\
13 & 13 & 7 & 0.453 & 0.625 & 0.724 & 0.271 & 0.099 \\
3 & 03 & 79 & 0.544 & 0.672 & 0.705 & 0.161 & 0.034 \\
4 & 04 & 29 & 0.464 & 0.663 & 0.705 & 0.242 & 0.042 \\
7 & 07 & 29 & 0.346 & 0.640 & 0.688 & 0.343 & 0.049 \\
2 & 02 & 69 & 0.549 & 0.673 & 0.684 & 0.135 & 0.011 \\
8 & 08 & 7 & 0.193 & 0.617 & 0.647 & 0.455 & 0.030 \\
11 & 11 & 22 & 0.189 & 0.564 & 0.615 & 0.426 & 0.051 \\
10 & 10 & 49 & 0.302 & 0.559 & 0.612 & 0.310 & 0.052 \\
12 & 12 & 32 & 0.201 & 0.560 & 0.586 & 0.386 & 0.027 \\
15 & 15 & 14 & 0.183 & 0.392 & 0.476 & 0.293 & 0.084 \\
9 & 09 & 2 & 0.077 & 0.312 & 0.412 & 0.335 & 0.100 \\
\end{longtable}
}

\begin{center}\rule{0.5\linewidth}{0.5pt}\end{center}

\subsection{Synthesis: What the Results
Say}\label{synthesis-what-the-results-say}

\begin{enumerate}
\def\labelenumi{\arabic{enumi}.}
\tightlist
\item
  FT \textgreater{} Local is universal at group level (18/18 groups).
\item
  FT \textgreater{} Global is also universal at group level (18/18
  groups).
\item
  Basin-level distributions show broad positive shift, not isolated
  outliers.
\item
  Transfer gains are strongest in lower-data or lower-baseline groups.
\end{enumerate}

\begin{center}\rule{0.5\linewidth}{0.5pt}\end{center}

\subsection{Practical Implications for
Deployment}\label{practical-implications-for-deployment}

\begin{itemize}
\tightlist
\item
  For new target regions with limited labels, start from a global model.
\item
  Use local fine-tuning instead of local-from-scratch as default.
\item
  Keep global-only evaluation as a diagnostic baseline.
\item
  Monitor groups with weak gains for feature or split-quality issues.
\end{itemize}

\begin{center}\rule{0.5\linewidth}{0.5pt}\end{center}

\subsection{Limitations of Current RQ3
Study}\label{limitations-of-current-rq3-study}

\begin{itemize}
\tightlist
\item
  Single random seed (1001) in current report.
\item
  Limited global pretraining budget.
\item
  Metrics centered on NSE/KGE; signature-based diagnostics not yet
  added.
\item
  No SHAP attribution in this version.
\end{itemize}

\begin{center}\rule{0.5\linewidth}{0.5pt}\end{center}

\subsection{Next Steps (Already
Planned)}\label{next-steps-already-planned}

\begin{enumerate}
\def\labelenumi{\arabic{enumi}.}
\tightlist
\item
  Multi-seed reruns with confidence intervals.
\item
  Longer global pretraining budget for stronger base weights.
\item
  SHAP-based transferability driver analysis.
\item
  Error decomposition by flow regime and event intensity.
\end{enumerate}

\begin{center}\rule{0.5\linewidth}{0.5pt}\end{center}

\subsection{Project Conclusion (RQ3)}\label{project-conclusion-rq3}

\begin{enumerate}
\def\labelenumi{\arabic{enumi}.}
\tightlist
\item
  This seminar project now provides a complete answer to RQ3.
\item
  Across 18 groups and 533 basins, \textbf{global + fine-tune}
  consistently outperforms local-from-scratch.
\item
  The transfer-learning strategy is therefore empirically supported and
  practically actionable.
\end{enumerate}

. . .

\begin{quote}
Fine-tuning a global hydrological model is a robust default strategy for
regional prediction under data constraints.
\end{quote}

\begin{center}\rule{0.5\linewidth}{0.5pt}\end{center}

\subsection{References}\label{references}

\protect\phantomsection\label{refs}
\begin{CSLReferences}{1}{1}
\bibitem[\citeproctext]{ref-addor2017}
Addor, Nans, Andrew J. Newman, Naoki Mizukami, and Martyn P. Clark.
2017. {``The {CAMELS} Data Set: Catchment Attributes and Meteorology for
Large-Sample Studies.''} \emph{Hydrology and Earth System Sciences} 21
(10): 5293--313. \url{https://doi.org/10.5194/hess-21-5293-2017}.

\bibitem[\citeproctext]{ref-cred2022}
Centre for Research on the Epidemiology of Disasters (CRED). 2022.
\emph{Disasters in Numbers 2022}. CRED.
\url{https://cred.be/sites/default/files/2022_EMDAT_report.pdf}.

\bibitem[\citeproctext]{ref-hochreiter1997}
Hochreiter, Sepp, and Juergen Schmidhuber. 1997. {``Long Short-Term
Memory.''} \emph{Neural Computation} 9 (8): 1735--80.
\url{https://doi.org/10.1162/neco.1997.9.8.1735}.

\bibitem[\citeproctext]{ref-hoedt2021}
Hoedt, Pieter-Jan, Frederik Kratzert, Daniel Klotz, et al. 2021.
{``{MC-LSTM}: Mass-Conserving {LSTM}.''} \emph{Proceedings of the 38th
International Conference on Machine Learning (ICML)}, 4275--86.
\url{https://arxiv.org/abs/2101.05186}.

\bibitem[\citeproctext]{ref-ipcc2021}
IPCC. 2021. {``Global {Carbon} and {Other} {Biogeochemical} {Cycles} and
{Feedbacks}.''} In \emph{Climate {Change} 2021 -- {The} {Physical}
{Science} {Basis}: {Working} {Group} {I} {Contribution} to the {Sixth}
{Assessment} {Report} of the {Intergovernmental} {Panel} on {Climate}
{Change}}. Cambridge University Press.

\bibitem[\citeproctext]{ref-kratzert2022joss}
Kratzert, Frederik, Martin Gauch, Grey Nearing, and Daniel Klotz. 2022.
{``{NeuralHydrology} --- {A} {Python} Library for Deep Learning Research
in Hydrology.''} \emph{Journal of Open Source Software} 7 (71): 4050.
\url{https://doi.org/10.21105/joss.04050}.

\bibitem[\citeproctext]{ref-kratzert2018}
Kratzert, Frederik, Daniel Klotz, Claire Brenner, Karsten Schulz, and
Mathew Herrnegger. 2018. {``Rainfall--Runoff Modelling Using Long
Short-Term Memory ({LSTM}) Based Neural Networks.''} \emph{Hydrology and
Earth System Sciences} 22 (11): 6005--22.
\url{https://doi.org/10.5194/hess-22-6005-2018}.

\bibitem[\citeproctext]{ref-kratzert2019}
Kratzert, Frederik, Daniel Klotz, Mathew Herrnegger, Alden K. Sampson,
Sepp Hochreiter, and Grey S. Nearing. 2019. {``Towards Learning
Universal, Regional, and Local Hydrological Behaviors via Machine
Learning Applied to Large-Sample Datasets.''} \emph{Hydrology and Earth
System Sciences} 23 (12): 5089--110.
\url{https://doi.org/10.5194/hess-23-5089-2019}.

\bibitem[\citeproctext]{ref-kratzert2023caravan}
Kratzert, Frederik, Grey Nearing, Jonathan Frame, et al. 2023.
{``Caravan -- {A} Global Community Dataset for Large-Sample
Hydrology.''} \emph{Scientific Data} 10 (1): 61.
\url{https://doi.org/10.1038/s41597-023-01975-w}.

\bibitem[\citeproctext]{ref-milly2008}
Milly, P. C. D., Julio Betancourt, Malin Falkenmark, et al. 2008.
{``Stationarity Is Dead: Whither Water Management?''} \emph{Science} 319
(5863): 573--74. \url{https://doi.org/10.1126/science.1151915}.

\end{CSLReferences}




\end{document}
